\section*{Capítulo \ref{ch:b3:molecular-evolution} --- Evolución molecular y filogenómica}

\begin{solution}{exo:b3:molecular-evolution:1}
$k = \mu$: el ritmo neutro de sustitución por sitio y generación es
igual al ritmo de mutación. En cada generación surgen $2N\mu$ mutaciones
neutras nuevas y cada una tiene una probabilidad $1/2N$ de fijarse; una
población mayor produce más mutaciones y da a cada una una probabilidad
proporcionalmente menor, y los dos efectos se cancelan exactamente.
\end{solution}

\begin{solution}{exo:b3:molecular-evolution:2}
\omterm{def:b3:genomics:comparative}{Ortólogos}: genes de dos especies descendientes de un solo gen de su
antepasado común ---la $\beta$-globina humana y la $\beta$-globina de
ratón. \omterm{def:b3:genomics:comparative}{Parálogos}: genes de un mismo genoma descendientes de una
duplicación ---la $\alpha$- y la $\beta$-globina humanas, o la $\beta$ y
la $\gamma$. \omterm{def:b3:molecular-evolution:duplication}{Seudogén}: una copia decaída que ya no codifica una proteína
---el gen $\psi\beta$ del grupo de la $\beta$-globina humana, con sus
codones de parada incluidos.
\end{solution}

\begin{solution}{exo:b3:molecular-evolution:3}
$\omega$ compara el \omterm{thm:b3:molecular-evolution:neutral}{ritmo de sustitución} que cambia aminoácidos con el
\omterm{thm:b3:molecular-evolution:neutral}{ritmo de sustitución} silenciosa, que es la línea de base neutra.
$0.02$: \omterm{def:b3:molecular-evolution:dnds}{selección purificadora} fuerte, con el \qty{98}{\%} de los
cambios de aminoácido eliminados ---actina, \omterm{def:b3:chromatin-epigenetics:nucleosome}{histona} H3. $1.0$: sin
restricción ---un \omterm{def:b3:molecular-evolution:duplication}{seudogén}. $2.5$: \omterm{def:b3:molecular-evolution:dnds}{selección positiva}, con los cambios
de aminoácido favorecidos ---el surco de unión a péptidos del MHC, la
envuelta del VIH, la lisozima de los rumiantes.
\end{solution}

\begin{solution}{exo:b3:molecular-evolution:4}
Las proteínas comparadas entre mamíferos cambian a alrededor de una
sustitución por sitio y por mil millones de años; sobre un genoma de
mil millones de sitios codificantes y una generación de unos pocos
años, eso es del orden de una sustitución fijada por generación, o una
cada dos años. El coste de la selección de Haldane permite alrededor de
una sustitución seleccionada cada trescientas generaciones. Los ritmos
difieren en un factor de cien o más, de modo que casi todos los cambios
fijados han de ser neutros.
\end{solution}

\begin{solution}{exo:b3:molecular-evolution:5}
Neutra: $1/2N = 10^{-4}$. $s = 0.001$, con $4Ns = 20$: $(1 - \mathrm{e}^{-0.002})/
(1 - \mathrm{e}^{-20}) = 0.002$. $s = 0.01$: $0.0198$. $s = -0.001$: $(1 -
\mathrm{e}^{0.002})/(1 - \mathrm{e}^{20}) = 0.002\,\mathrm{e}^{-20} =
4\times 10^{-12}$. Ninguna es efectivamente neutra: para eso haría falta
$|4Ns|
\lesssim 1$, es decir, $|s| \lesssim 5\times 10^{-5}$; un coeficiente de
selección de una décima de uno por ciento ya es decisivo en una
población de cinco mil.
\end{solution}

\begin{solution}{exo:b3:molecular-evolution:6}
Sin corregir: $T = 0.12/(2\times 2\times 10^{-9}) = \qty{30}{Myr}$. Con
Jukes--Cantor: $d = -\tfrac{3}{4}\ln(1 - 0.16) = 0.131$, $T =
\qty{33}{Myr}$. La corrección alcanza un factor de dos cuando $-\tfrac{3}{4}
\ln(1 - 4p/3) = 2p$, en $p \approx 0.6$: para entonces la diferencia
bruta ha perdido casi toda su información.
\end{solution}

\begin{solution}{exo:b3:molecular-evolution:7}
$p_{N} = 7/700 = 0.010$, $p_{S} = 10/200 = 0.050$, $\omega = 0.2$:
\omterm{def:b3:molecular-evolution:dnds}{selección purificadora} que elimina cuatro quintas partes de los cambios
de aminoácido. Con $35$ y $10$: $p_{N} = 0.05$, $\omega = 1$: sin
restricción detectable ---un \omterm{def:b3:molecular-evolution:duplication}{seudogén}, o un gen en el que la \omterm{def:b3:molecular-evolution:dnds}{selección
positiva} en algunos sitios equilibra la purificadora en otros.
\end{solution}

\begin{solution}{exo:b3:molecular-evolution:8}
$k = d/2T = 0.012/(1.3\times 10^{7}) = 9.2\times 10^{-10}$ por sitio y
año; $\mu = 25k = 2.3\times 10^{-8}$ por sitio y generación; $6\times
10^{9}\times 2.3\times 10^{-8} \approx 140$ mutaciones nuevas por genoma
diploide y generación. La secuenciación directa de padres e hijos
encuentra unas setenta: la estimación por divergencia queda inflada por
el polimorfismo ya presente en la población ancestral, que añade unos
cientos de miles de generaciones a la separación aparente.
\end{solution}

\begin{solution}{exo:b3:molecular-evolution:9}
$\mathrm{e}^{-T/2N} = \mathrm{e}^{-1} = 0.37$; fracción discordante, $\tfrac{2}{3}
\times 0.37 = 0.25$. Para el \qty{10}{\%}: $\mathrm{e}^{-T/2N} = 0.15$, $T =
2N\ln(1/0.15) = \num{190000}$ generaciones.
\end{solution}

\begin{solution}{exo:b3:molecular-evolution:10}
Cuando $s \to 0$: $1 - \mathrm{e}^{-2s} \approx 2s$ y $1 - \mathrm{e}^{-4Ns}
\approx 4Ns$, y la razón vale $1/2N$. Para $4Ns \gg 1$ el denominador es
$1$ y el numerador $\approx 2s$. Para $s < 0$ con $4N|s| \gg 1$:
numerador $1 -
\mathrm{e}^{2|s|} \approx -2|s|$, denominador $1 - \mathrm{e}^{4N|s|} \approx
-\mathrm{e}^{4N|s|}$, de modo que $P \approx 2|s|\,\mathrm{e}^{-4N|s|}$. Cien
veces por debajo de la neutra con $N = 10^{4}$: $2|s|\,\mathrm{e}^{-4\times 10^{4}|s|} =
5\times 10^{-7}$; $|s| = 1.6\times 10^{-4}$ da $3.2\times 10^{-4}\times
\mathrm{e}^{-6.4} = 5.3\times 10^{-7}$. Así que $|s| \approx 1.6\times 10^{-4}$
y $4N|s| \approx 6.5$: una desventaja de una centésima de uno por
ciento basta, en diez mil individuos, para que la fijación sea cien
veces más rara que por azar.
\end{solution}

\begin{solution}{exo:b3:molecular-evolution:11}
Con el punto de separación a una distancia $a$ de $O$ por ambos
caminos, $a + m =
0.30$, $a + h = 0.24$ y $m + h = 0.20$: $m - h = 0.06$, de modo que
$m = 0.13$ y $h
= 0.07$, con una razón de $1.9$. Un reloj estricto por generación daría
la razón de los recuentos de generaciones, $25/0.5 = 50$. El $1.9$
observado significa que el ratón ha acumulado solo el doble de cambio
anual que el ser humano pese a tener cincuenta veces más generaciones:
el ritmo de mutación por generación ha de ser unas $25$ veces mayor en
el ser humano ---más divisiones celulares de la línea germinal por
generación--- de modo que por año los dos linajes difieran solo en un
factor de dos.
\end{solution}

\begin{solution}{exo:b3:molecular-evolution:12}
Las sustituciones que crean codones de parada llegan a $1000\times 2\times 10^{-9}\times
0.04 = 8\times 10^{-8}$ por año: cabe esperar la primera al cabo de unos
\qty{12}{Myr} (los desplazamientos de pauta reducen esto casi a la
mitad). En esa ventana, una mutación beneficiosa, una de cada $10^{3}$ a
$10^{5}$ de todas las mutaciones, tiene mucha menos probabilidad de
llegar que un codón de parada, e incluso cuando llega se fija con
probabilidad solo $2s$. De modo que la copia suele estar muerta antes de
que la selección pueda hallarle un uso ---la semivida observada de los
duplicados en los vertebrados es de unos pocos millones de años.
\end{solution}

\begin{solution}{pb:b3:molecular-evolution:1}
\textbf{1.} $k = 0.012/(2\times 6.5\times 10^{6}) = 9.2\times 10^{-10}$ por
sitio y año; $\mu = 25k = 2.3\times 10^{-8}$ por sitio y generación.
\textbf{2.} $6\times 10^{9}\times 2.3\times 10^{-8} \approx 140$
mutaciones nuevas por genoma diploide; el \qty{1.5}{\%} de ellas, unas
$2$, en secuencia codificante.
\textbf{3.} Por genoma haploide, $3\times 10^{9}\times 2.3\times 10^{-8}
\approx 70$ sustituciones fijadas por generación, frente al $1/300$ de
Haldane: veinte mil veces más de lo que la selección podría impulsar, de
modo que la abrumadora mayoría son neutras.
\textbf{4.} $2N\mu = 10^{5}\times 2.3\times 10^{-8} = 2.3\times 10^{-3}$
mutaciones nuevas por sitio y generación en la población; se fijará
$1/2N =
10^{-5}$ de ellas.
\textbf{5.} $4N = \num{200000}$ generaciones, \qty{5}{Myr} ---casi tanto
como el tiempo transcurrido desde la separación.
\textbf{6.} La divergencia cuenta las mutaciones que surgieron a lo
largo de los dos linajes separados desde su antepasado común; cada una
se fija dentro de su propio linaje, y el ritmo a largo plazo es $\mu$
sea cual sea el tiempo que tarde cada una, ya que las mutaciones surgen
de forma continua y el retardo solo demora el flujo, sin cambiarlo. (La
única corrección es por el polimorfismo presente en la separación, que
hace que el antepasado común de dos alelos sea más antiguo que la
separación de las especies.)
\textbf{7.} Neutra, $10^{-5}$; $s = 0.001$: $0.002$; $s = 0.01$: $0.0198$;
$s = 0.1$: $1 - \mathrm{e}^{-0.2} = 0.18$.
\textbf{8.} $s = 1/4N = 5\times 10^{-6}$: por debajo de esto la deriva
anega a la selección y la mutación se comporta como neutra; por encima,
la selección gobierna las probabilidades.
\textbf{9.} $s = -0.001$: $P = 0.002\,\mathrm{e}^{-200}$, cero a todos los
efectos ---$10^{-85}$ de la neutra. $s = -10^{-5}$, con $4N|s| = 2$: $P =
2\times 10^{-5}/(\mathrm{e}^{2} - 1) = 3.1\times 10^{-6}$, el \qty{31}{\%}
de la neutra: las mutaciones levemente deletéreas sí se fijan.
\textbf{10.} Sustituciones beneficiosas por sitio y generación: $2N\times
10^{-5}\mu\times 2s = 10^{5}\times 10^{-5}\times 0.02\,\mu = 0.02\mu$;
neutras: $\approx\mu$. Fracción adaptativa,
$0.02/1.02 \approx \qty{2}{\%}$. Aunque cada mutación beneficiosa tiene
dos mil veces más probabilidad de fijarse que una neutra, son tan raras
que solo una sustitución de cada cincuenta es adaptativa: la selección
actúa y apenas se asoma.
\textbf{11.} $p_{N} = 46/920 = 0.050$, $p_{S} = 84/280 = 0.30$.
Corregidos: $d_{N} = -\tfrac{3}{4}\ln(1 - 0.0667) = 0.052$, $d_{S} = -\tfrac{3}{4}
\ln(1 - 0.40) = 0.38$. $\omega = 0.14$.
\textbf{12.} \omterm{def:b3:molecular-evolution:dnds}{Selección purificadora} que elimina alrededor del
\qty{86}{\%} de los cambios de aminoácido ---un gen típico. $k = d_{S}/2T = 0.38/(1.8\times
10^{8}) = 2.1\times 10^{-9}$ por sitio y año: el doble del valor
humano--chimpancé, como cabe esperar de una comparación que incluye el
linaje de tictac rápido de los roedores; el mismo orden de magnitud.
\textbf{13.} $d = -\ln(1 - 0.55) = 0.80$ por sitio; $T = 0.80/(2\times
10^{-9}) = \qty{400}{Myr}$.
\textbf{14.} Dos cadenas distintas forman el tetrámero
$\alpha_{2}\beta_{2}$, cuya unión cooperativa del oxígeno, \omterm{prop:b3:structural-biology:mechanism}{efecto Bohr} y
control alostérico por el 2,3-bisfosfoglicerato dependen de la
interfaz entre subunidades distintas; y el grupo $\beta$ pudo
diversificarse después en cadenas embrionarias, fetales y adultas con
afinidades distintas.
\textbf{15.} $2kT = 2\times 8\times 10^{-9}\times 4\times 10^{7} = 0.64$
sustituciones por sitio; con la saturación entre veinte aminoácidos, la
diferencia bruta es de alrededor de
$0.95(1 - \mathrm{e}^{-0.64/0.95}) \approx 0.47$. A \qty{500}{Myr},
$d = 8$: todos los sitios han cambiado muchas veces y la diferencia
bruta se sitúa en su techo del \qty{95}{\%}, sin llevar información
alguna sobre el tiempo.
\textbf{16.} $10^{-11}\times 10^{9}\times 100\times 2 = 2$ sustituciones
entre dos linajes en 100 sitios a lo largo de mil millones de años. Para
una separación de \qty{10}{Myr} la esperanza es $0.02$: casi siempre
cero diferencias, sin resolución.
\textbf{17.} $1200\times 2\times 10^{-9}\times 0.04 = 9.6\times 10^{-8}$
por año: cabe esperar la primera parada al cabo de unos \qty{10}{Myr}.
\textbf{18.} Duplicar el genoma entero duplica todos los genes a la vez,
de modo que se conserva la estequiometría de las proteínas que
interaccionan y ningún desequilibrio de dosis selecciona contra las
copias; una duplicación en tándem aislada cambia la dosis de un solo
gen y a menudo es deletérea, o bien es redundante de inmediato y queda
libre para decaer.
\textbf{19.} $\mathrm{e}^{-\num{80000}/\num{100000}} = \mathrm{e}^{-0.8} = 0.45$.
\textbf{20.} $\tfrac{2}{3}\times 0.45 = 0.30$: coherente con el
\qty{30}{\%} observado.
\textbf{21.} La mitad cada una: el \qty{15}{\%} de todos los árboles
agrupa al humano con el gorila y el \qty{15}{\%} al chimpancé con el
gorila. Un exceso claro de uno indicaría flujo génico entre esas dos
especies después de su separación, cosa que la deriva no puede
producir.
\textbf{22.} $N = \num{10000}$: $\mathrm{e}^{-4} = 0.018$, y una
discordancia del \qty{1.2}{\%}. El \qty{30}{\%} observado exige, por
tanto, una población ancestral de unos cincuenta mil ---varias veces el
tamaño efectivo del ser humano actual.
\textbf{23.} La concatenación supone que todo gen tiene el \omterm{prop:b3:molecular-evolution:ils}{árbol de
especies}; con \omterm{prop:b3:molecular-evolution:ils}{clasificación incompleta de linajes} los genes tienen
muchos árboles, y para ramas internas cortas el \omterm{prop:b3:molecular-evolution:ils}{árbol de genes} más
común puede diferir del \omterm{prop:b3:molecular-evolution:ils}{árbol de especies}, de modo que más datos hacen
más segura la respuesta equivocada. Un método coalescente calcula, para
cada \omterm{prop:b3:molecular-evolution:ils}{árbol de especies} candidato, la distribución esperada de \omterm{prop:b3:molecular-evolution:ils}{árboles de
genes}, y elige el \omterm{prop:b3:molecular-evolution:ils}{árbol de especies} que mejor explica las frecuencias
observadas.
\textbf{24.} La \omterm{prop:b3:molecular-evolution:ils}{clasificación incompleta de linajes} ocurrió en la
población ancestral común de todos los humanos modernos, de modo que
haría a toda población humana igualmente emparentada con los
neandertales. Un exceso en los no africanos significa que los alelos
neandertales entraron después de que los antepasados de los no
africanos salieran de África: mestizaje, hace unos cincuenta mil años.
\textbf{25.} $\mu \approx 2.3\times 10^{-8}$ por sitio y generación;
$\omega \approx 0.14$; alrededor del \qty{30}{\%} de los \omterm{prop:b3:molecular-evolution:ils}{árboles de
genes} discordantes con el \omterm{prop:b3:molecular-evolution:ils}{árbol de especies}.
\end{solution}
